
\subsubsection{Weight-Space Learning}

Unterthiner et al. (2020) established the weight-to-accuracy prediction task, demonstrating $R^2$ > 0.98 using per-layer statistics (mean, standard deviation, spectral norm) on 120K CNN models. This work proves the task is solvable but sidesteps the symmetry problem through feature engineering. Eilertsen et al. (2020) extended this to ''classifying the classifier,'' predicting training hyperparameters from weight distributions.

Schürholt et al. (2022) released Model Zoos, a benchmark of 50K+ neural network checkpoints enabling systematic weight-space research. Their follow-up work, SANE (Schürholt et al., 2024), introduced scalable weight embeddings via sequential token processing.

\subsubsection{Permutation Symmetry in Neural Networks}

Ainsworth et al. (2023) formalized permutation symmetry in weight space through Git Re-Basin, showing that independently trained networks can be merged by aligning their hidden unit orderings. Sharma et al. (2024) analyzed variance collapse in model merging, demonstrating that permutation alignment is necessary for meaningful interpolation.

\subsubsection{Permutation-Equivariant Architectures}

Zhou et al. (2023) introduced Neural Functional Networks (NFN), providing permutation-equivariant layers for processing neural network weights. Zaheer et al. (2017) established DeepSets, the foundational framework for permutation-invariant/equivariant set functions. Tran-Viet et al. (2024) extended NFN to transformer architectures and released a dataset of 125K transformer checkpoints.

Prior work either uses hand-crafted invariant features or demonstrates equivariant architectures without systematic baseline comparison. The present work provides this comparison, finding that equivariance is not merely efficient but necessary for learning from raw weights.
